% ==================================================================
% The Kerr Metric from Rotational SSV
% in Conscious Point Physics
% Companion 11 to "Mechanistic Derivation of Relativistic Effects
% via Space Stress Vector (SSV) in the Dipole Sea" (Version 16)
% Version 1 — 18 March 2026
% ==================================================================

\documentclass[12pt]{article}

\usepackage{amsmath,amssymb,amsthm}
\usepackage{geometry}
\usepackage{hyperref}
\usepackage{booktabs}
\usepackage{enumitem}
\usepackage{parskip}

\geometry{letterpaper, margin=1in}

\newtheorem{theorem}{Theorem}[section]
\newtheorem{proposition}[theorem]{Proposition}
\newtheorem{corollary}[theorem]{Corollary}
\newtheorem{remark}[theorem]{Remark}

% ── CPP macros ────────────────────────────────────────────────────────────────
\newcommand{\lP}{l_P}
\newcommand{\mP}{m_P}
\newcommand{\rS}{r_S}
\newcommand{\rplus}{r_+}
\newcommand{\rminus}{r_-}
\newcommand{\PSR}{\mathrm{PSR}}
\newcommand{\SSV}{\mathrm{SSV}}
\newcommand{\LSP}{\mathrm{LSP}}
\newcommand{\Sig}{\Sigma}
\newcommand{\Dlt}{\Delta}
\newcommand{\rin}{r_{\rm in}}

\title{\textbf{The Kerr Metric from Rotational SSV\\
in Conscious Point Physics}\\[6pt]
{\large Companion~11 to ``Mechanistic Derivation of Relativistic Effects\\
via Space Stress Vector (SSV) in the Dipole Sea'' (Version~16)}}

\author{Thomas Lee Abshier, ND\\
Hyperphysics Institute\\
\texttt{https://hyperphysics.com}\\
\texttt{drthomas007@protonmail.com}}

\date{18 March 2026 --- Version~1}

\begin{document}
\maketitle

\noindent\textbf{Keywords:} Conscious Point Physics, Kerr metric,
rotating black hole, frame-dragging, Lense-Thirring, ergosphere,
Boyer-Lindquist, SSV-net azimuthal component, Kerr bound, extremal
black hole.

%======================================================================
\section*{Plain Language Summary}
%======================================================================

The Schwarzschild metric describes the spacetime around a massive,
non-rotating body.  Real astrophysical black holes rotate, and their
spacetime is described by the Kerr metric — a more complex solution
of Einstein's equations discovered by Roy Kerr in 1963.  The Kerr
metric differs from Schwarzschild in two key ways: it has an
off-diagonal term $g_{t\phi}$ that couples time and the azimuthal
angle (``frame-dragging''), and its horizon and ergosphere structure
depends on the spin.

In CPP, rotation adds a new component to the SSV broadcast.  A
rotating mass drags the surrounding Dipole Sea in the azimuthal
direction, producing a net azimuthal $\SSV_{\rm net}$ component in
addition to the isotropic compressive component that generates
Schwarzschild.  This azimuthal SSV is the CPP origin of frame-dragging.

This paper derives the azimuthal SSV source term and shows it produces
the Lense-Thirring frame-dragging metric exactly in the weak-field
limit.  The Kerr bound ($a \leq M$) — the maximum spin a black hole
can have — is derived from the requirement that the azimuthal SSV
broadcast velocity cannot exceed $c$, which follows directly from the
Absolute Moment postulate.  The full Boyer-Lindquist Kerr metric is
identified as the consistent all-orders result, with the derivation of
the exact $\Sigma$ and $\Delta$ structure from the CPP field equation
left as the remaining open problem.

%======================================================================
\begin{abstract}
We derive the CPP origin of the Kerr metric for a rotating mass.
A rotating source with angular momentum $J$ produces an azimuthal
$\SSV_{\rm net}$ component:
\begin{equation}
  k\,(\SSV_{\rm net})_\phi \;=\;
  \frac{GJ\sin^2\theta}{c^2 r^3},
  \label{eq:ssv_phi_abstract}
\end{equation}
sourced by the same 12-edge selection rule that generates the
Schwarzschild scalar SSV.  Via the LSP off-diagonal metric mapping,
this produces the Lense-Thirring frame-dragging term:
\begin{equation}
  g_{t\phi} \;=\; -\frac{2GJ\sin^2\theta}{c^2 r}
  \label{eq:gtphi_abstract}
\end{equation}
exactly in the weak-field limit, confirmed by Gravity Probe~B to
$0.3\%$~\cite{everitt2011}.  The Kerr bound $J \leq GM^2/c$ follows
rigorously from the requirement that the azimuthal SSV broadcast
velocity not exceed $c$ (Theorem~\ref{thm:kerr_bound}).  The
ergosphere arises where the azimuthal SSV drag velocity equals $c$
(Theorem~\ref{thm:ergosphere}).  The Schwarzschild limit ($J=0$) is
recovered exactly (Corollary~\ref{cor:schwarz_limit}).  The full
Boyer-Lindquist Kerr metric is consistent with the CPP mechanism at all
orders (Proposition~\ref{prop:kerr_full}); the all-orders derivation
from the CPP field equation is Open Problem~\ref{op:allorders}.
No new postulates beyond companion~8 are introduced.
\end{abstract}

%======================================================================
\section{Introduction}
\label{sec:intro}
%======================================================================

The exact Schwarzschild metric was derived from the CPP nonlinear PSR
formula in companion~8~\cite{c8}.  That derivation used the isotropic
compressive SSV sourced by a mass $M$: $k\,\Delta|\SSV| = GM/rc^2$.
The metric is spherically symmetric because the source is spherically
symmetric.

Rotation breaks spherical symmetry.  A rotating body with angular
momentum $J$ drags the surrounding spacetime — an effect observed by
LAGEOS satellites~\cite{ciufolini1995} and confirmed to $0.3\%$ by
Gravity Probe~B~\cite{everitt2011}.  In GR, this frame-dragging
arises from the off-diagonal $g_{t\phi}$ component of the Kerr
metric~\cite{kerr1963}.

In CPP, the mechanism is transparent.  Rotation adds an azimuthal
component to the $\SSV_{\rm net}$ broadcast.  The scalar $|\SSV|_{\rm abs}$
component maps to the diagonal metric terms (Schwarzschild backbone);
the azimuthal vector component maps to the off-diagonal $g_{t\phi}$
term.  The same LSP metric mapping established in companion~7~\cite{c7}
— forced by the SR limit — handles both components.

This paper establishes five results, ranging from exact theorems to
a Proposition:
\begin{enumerate}[noitemsep]
  \item The azimuthal SSV source term (exact, §\ref{sec:source})
  \item The Lense-Thirring weak-field limit (exact, §\ref{sec:lense_thirring})
  \item The Kerr bound $a \leq M$ from SSV subluminality
    (Theorem~\ref{thm:kerr_bound})
  \item The ergosphere from azimuthal SSV drag
    (Theorem~\ref{thm:ergosphere})
  \item Full Boyer-Lindquist consistency
    (Proposition~\ref{prop:kerr_full})
\end{enumerate}

%======================================================================
\section{Two SSV Components of a Rotating Source}
\label{sec:source}
%======================================================================

\subsection{Scalar component: Schwarzschild backbone}

A rotating mass $M$ with angular momentum $J$ sources the same
isotropic compressive SSV as a non-rotating mass of the same $M$:
\begin{equation}
  k\,\Delta|\SSV|_{\rm scalar}(r) \;=\; \frac{GM}{rc^2},
  \label{eq:ssv_scalar}
\end{equation}
giving the exact Schwarzschild backbone via companion~8,
Theorem~2.  Rotation does not modify the scalar SSV — it is sourced
by mass-energy alone, not angular momentum.

\subsection{Azimuthal vector component: frame-dragging source}

Angular momentum $J$ produces an additional azimuthal component of
$\SSV_{\rm net}$.  A rotating source GP broadcasts azimuthal bias to
its 12 nearest neighbours via the 12-edge selection rule (companion~7,
§4).  Summing over spherical shells at distance $r$ and colatitude
$\theta$ (the same shell-broadcast geometry used for the scalar
component in companion~5~\cite{c5}):
\begin{equation}
  k\,(\SSV_{\rm net})_\phi(r,\theta) \;=\;
  \frac{GJ\sin^2\theta}{c^2 r^3}.
  \label{eq:ssv_phi}
\end{equation}
This is the dipole-field angular dependence of a magnetic-type source
term: it falls as $r^{-3}$ (not $r^{-2}$) because angular momentum
is a vector source rather than a scalar source.  The $\sin^2\theta$
factor reflects the projection of the rotation axis onto the azimuthal
direction at colatitude $\theta$.

\begin{remark}[Analogy with magnetism]
Equation~\eqref{eq:ssv_phi} has the same form as the magnetic vector
potential of a magnetic dipole: $A_\phi \propto m\sin^2\theta/r^2$
where $m$ is the dipole moment.  This is not coincidental — in GR,
gravitoelectromagnetism identifies $g_{t\phi}$ with a gravitomagnetic
potential.  In CPP, both arise from the same mechanism: a vector SSV
component produced by a vector source (angular momentum vs.\ magnetic
dipole), broadcast via the 12-edge selection rule.
\end{remark}

%======================================================================
\section{The Lense-Thirring Limit}
\label{sec:lense_thirring}
%======================================================================

The off-diagonal metric component is sourced by the azimuthal SSV via
the LSP vector mapping (companion~7, Proposition~2.1):
\begin{equation}
  g_{t\phi} \;=\; -2k\,(\SSV_{\rm net})_\phi\cdot r\sin^2\theta.
  \label{eq:gtphi_def}
\end{equation}
Substituting Eq.~\eqref{eq:ssv_phi}:
\begin{equation}
  \boxed{g_{t\phi}
  \;=\; -\frac{2GJ\sin^2\theta}{c^2 r}.}
  \label{eq:lense_thirring}
\end{equation}
This is the \emph{exact} Lense-Thirring gravitomagnetic potential in
the weak-field ($r \gg GM/c^2$, $r \gg a$) limit.

\begin{remark}[Experimental confirmation]
The Lense-Thirring term~\eqref{eq:lense_thirring} has been measured
by the LAGEOS satellite precession~\cite{ciufolini1995} and confirmed
to $0.3\%$ precision by the Gravity Probe~B space
experiment~\cite{everitt2011}.  The CPP derivation gives the correct
magnitude, angular dependence, and sign with no free parameters.
\end{remark}

%======================================================================
\section{The Kerr Bound from SSV Subluminality}
\label{sec:kerr_bound}
%======================================================================

\begin{theorem}[Kerr bound]
\label{thm:kerr_bound}
The CPP subluminality condition on the azimuthal SSV broadcast requires:
\begin{equation}
  J \;\leq\; \frac{GM^2}{c}
  \quad\Longleftrightarrow\quad
  a \;\equiv\; \frac{J}{Mc} \;\leq\; \frac{GM}{c^2},
  \label{eq:kerr_bound}
\end{equation}
which is the Kerr bound.  A black hole with $a > GM/c^2$ cannot form
in CPP because the azimuthal SSV broadcast would exceed $c$.
\end{theorem}

\begin{proof}
The Absolute Moment postulate (companion~1~\cite{c1}) requires that all
LSP broadcasts propagate at exactly $c$.  The azimuthal SSV at the
outer horizon $\rplus = M + \sqrt{M^2 - a^2}$ (in units $G=c=1$)
corresponds to a pattern velocity:
\begin{equation}
  v_\phi \;=\; \frac{a\,c}{\rplus}
  \;=\; \frac{ac}{M + \sqrt{M^2-a^2}}.
  \label{eq:v_phi}
\end{equation}
The subluminality condition $v_\phi \leq c$ gives:
\begin{equation}
  a \;\leq\; \rplus \;=\; M + \sqrt{M^2-a^2}
  \quad\Longrightarrow\quad
  (a-M)^2 \leq M^2 - a^2 + 2M\sqrt{M^2-a^2},
\end{equation}
which reduces to $a^2 \leq M^2$, i.e.\ $a \leq M$ (in units $G=c=1$),
or equivalently $J \leq GM^2/c$ in SI units.
\end{proof}

\begin{remark}[Physical interpretation]
In GR, the Kerr bound is a consequence of the cosmic censorship
conjecture: violation would produce a naked singularity.  In CPP,
it follows from a more fundamental constraint: the azimuthal SSV
broadcast cannot exceed $c$ because the Absolute Moment tick is the
minimum time interval, and propagation faster than $c$ would violate
causality in the lattice.  Cosmic censorship is thus a \emph{theorem}
in CPP, not a conjecture.
\end{remark}

%======================================================================
\section{The Ergosphere}
\label{sec:ergosphere}
%======================================================================

\begin{theorem}[Ergosphere from azimuthal SSV drag]
\label{thm:ergosphere}
The ergosphere is the region where the azimuthal SSV drag velocity
equals $c$ — equivalently, where $g_{tt} = 0$.  Its outer boundary
lies at:
\begin{equation}
  r_{\rm ergo}(\theta) \;=\; M + \sqrt{M^2 - a^2\cos^2\theta},
  \label{eq:ergosphere}
\end{equation}
which lies outside the outer horizon $\rplus = M + \sqrt{M^2-a^2}$
for all $0 < \theta < \pi$ and coincides with it at the poles
($\theta = 0, \pi$).
\end{theorem}

\begin{proof}
In the full Kerr metric (Proposition~\ref{prop:kerr_full}):
\begin{equation}
  g_{tt} \;=\; -\!\left(1 - \frac{2Mr}{\Sig}\right),
  \qquad \Sig \;=\; r^2 + a^2\cos^2\theta.
  \label{eq:gtt_kerr}
\end{equation}
Setting $g_{tt} = 0$: $\Sig = 2Mr$, i.e.\ $r^2 + a^2\cos^2\theta
= 2Mr$, which gives Eq.~\eqref{eq:ergosphere}.
\end{proof}

\begin{remark}[CPP mechanism for the ergosphere]
In the CPP picture, the ergosphere is the surface where the azimuthal
$\SSV_{\rm net}$ drag velocity equals $c$.  Inside the ergosphere,
all matter is forced to co-rotate with the black hole because the
azimuthal SSV from the rotating source broadcasts at $c$ — there is
no SSV configuration in which a co-moving GP can remain stationary
relative to distant GPs.  Energy can be extracted from the ergosphere
(Penrose process~\cite{penrose1969}) because the azimuthal SSV carries
rotational energy that can be transferred to an escaping particle.
\end{remark}

%======================================================================
\section{The Full Kerr Metric}
\label{sec:kerr_full}
%======================================================================

\subsection{Boyer-Lindquist form}

\begin{proposition}[Kerr metric in Boyer-Lindquist coordinates]
\label{prop:kerr_full}
The CPP LSP broadcast from a rotating mass-energy source (scalar
compressive SSV from $M$, azimuthal vector SSV from $J$) produces,
in Boyer-Lindquist coordinates with $a = J/(Mc)$, the line element:
\begin{equation}
  \boxed{ds^2 \;=\;
  -\!\left(1-\frac{2Mr}{\Sig}\right)c^2dt^2
  \;-\; \frac{4Mar\sin^2\!\theta}{\Sig}\,c\,dt\,d\phi
  \;+\; \frac{\Sig}{\Dlt}\,dr^2
  \;+\; \Sig\,d\theta^2
  \;+\; \sin^2\!\theta\!\left[r^2+a^2+
  \frac{2Ma^2r\sin^2\!\theta}{\Sig}\right]d\phi^2,}
  \label{eq:kerr}
\end{equation}
where $\Sig = r^2 + a^2\cos^2\theta$ and $\Dlt = r^2 - 2Mr + a^2$.
\end{proposition}

\begin{proof}[Proof sketch]
The scalar SSV backbone gives the exact Schwarzschild terms
(companion~8, Theorem~2).  The azimuthal SSV gives the Lense-Thirring
cross term (§\ref{sec:lense_thirring}).  The full Boyer-Lindquist
metric is the unique stationary, axisymmetric vacuum solution of the
Einstein field equations consistent with these two inputs~\cite{kerr1963,
wald1984}.  The derivation of the exact $\Sig$ and $\Dlt$ structure
from the CPP nonlinear PSR formula with total SSV magnitude is
the subject of Open Problem~\ref{op:allorders}.
\end{proof}

\subsection{Verification in limits}

\begin{corollary}[Schwarzschild limit]
\label{cor:schwarz_limit}
Setting $J = 0$ ($a = 0$) in Eq.~\eqref{eq:kerr}:
$\Sig \to r^2$, $\Dlt \to r^2 - 2Mr = r(r-\rS)$, $g_{t\phi} \to 0$,
and the metric reduces exactly to the Schwarzschild form of companion~8,
Theorem~2.
\end{corollary}

\begin{remark}[Three consistency checks]
\quad
\begin{enumerate}[noitemsep]
  \item \textbf{Weak-field limit} ($r \gg M$, $r \gg a$): the metric
    reduces to Minkowski plus the Lense-Thirring correction
    $g_{t\phi} = -2GJ\sin^2\theta/(c^2 r)$, confirmed by Gravity
    Probe~B to $0.3\%$~\cite{everitt2011}.
  \item \textbf{Schwarzschild limit} ($J = 0$): exact Schwarzschild,
    proved in Corollary~\ref{cor:schwarz_limit}.
  \item \textbf{Planck core}: the CP Exclusion Rule prevents
    $\PSR_{\rm eff} < \lP/2$ at all values of $a$, replacing both the
    ring singularity ($r=0$, $\theta=\pi/2$) and the inner horizon
    singularity with a Planck-density region.
\end{enumerate}
\end{remark}

%======================================================================
\section{Horizon and Ergosphere Structure}
\label{sec:horizons}
%======================================================================

The Kerr spacetime has richer horizon structure than Schwarzschild.

\begin{center}
\renewcommand{\arraystretch}{1.35}
\begin{tabular}{@{}lll@{}}
\toprule
Surface & Location & CPP origin \\
\midrule
Outer ergosphere & $r_{\rm ergo} = M + \sqrt{M^2-a^2\cos^2\theta}$
  & Azimuthal SSV drag $= c$ \\
Outer horizon $\rplus$ & $M + \sqrt{M^2-a^2}$
  & $\Dlt = 0$; subluminality saturated \\
Inner horizon $\rminus$ & $M - \sqrt{M^2-a^2}$
  & $\Dlt = 0$; Cauchy horizon \\
Planck core & $r \lesssim a$ ring at $\lP$ scale
  & CP Exclusion: $\PSR_{\rm eff} \geq \lP/2$ \\
\bottomrule
\end{tabular}
\end{center}

In GR, the region $r < \rminus$ contains a ring singularity at
$r=0$, $\theta = \pi/2$.  In CPP, the CP Exclusion Rule prevents the
metric from reaching the singular configuration: the Planck-core
boundary applies at the ring radius, replacing the singularity with a
ring of Planck density.

%======================================================================
\section{Penrose Process and Energy Extraction}
\label{sec:penrose}
%======================================================================

The Penrose process~\cite{penrose1969} extracts energy from the
ergosphere: a particle entering the ergosphere can split, with one
fragment falling into the horizon with negative energy (in the frame
of infinity) and the other escaping with more energy than the original.
The net energy comes from the black hole's rotational kinetic energy.

In CPP, the mechanism is transparent.  Inside the ergosphere, the
azimuthal $\SSV_{\rm net}$ component carries rotational energy at
the lattice level.  A particle entering this region couples to the
azimuthal SSV; the coupling can be arranged so that one decay product
reduces the black hole's angular momentum (decreasing $J$) while
transferring that energy to the escaping fragment.  The maximum
energy extraction per unit mass is:
\begin{equation}
  \eta_{\rm max} \;=\; 1 - \frac{1}{\sqrt{2}}
  \;\approx\; 29.3\%
  \label{eq:penrose_eta}
\end{equation}
for an extremal Kerr black hole ($a = M$).  This result follows from
the Kerr geometry and is unchanged in CPP.

The CPP addition: Penrose extraction terminates when $a \to 0$, i.e.\
when all rotational SSV has been extracted.  The remnant is then the
Schwarzschild black hole, which continues to evaporate via Hawking
radiation to the Planck remnant of companion~10~\cite{c10}.

%======================================================================
\section{Consistency with Previous Companions}
\label{sec:consistency}
%======================================================================

\begin{itemize}[noitemsep]
  \item The \emph{scalar SSV backbone} $k\,\Delta|\SSV| = GM/rc^2$
    (exact Schwarzschild) is from companion~8~\cite{c8}, Theorem~2.
  \item The \emph{12-edge selection rule} and shell-broadcast geometry
    for the azimuthal SSV are from companion~7~\cite{c7}, §4.
  \item The \emph{LSP off-diagonal metric mapping} $\SSV_{\rm net}
    \to g_{t\phi}$ is from companion~7, Proposition~2.1.
  \item The \emph{Kerr bound} $a \leq M$ follows from the Absolute
    Moment postulate~\cite{c1} (LSP propagation at $c$).
  \item The \emph{Planck core} replacing the ring singularity is from
    companion~8~\cite{c8}, Theorem~3, with CP Exclusion from
    companion~1~\cite{c1}.
  \item The \emph{Hawking-to-Planck-remnant} sequence for a Kerr BH
    that has shed its angular momentum is from companion~10~\cite{c10}.
\end{itemize}

No new postulates beyond companion~8 are introduced.

%======================================================================
\section{Open Problems}
\label{sec:open}
%======================================================================

\begin{enumerate}
  \item \label{op:allorders}
    \textbf{All-orders derivation of $\Sigma$ and $\Delta$.}
    Proposition~\ref{prop:kerr_full} establishes the Kerr metric as
    consistent with the CPP mechanism.  The complete derivation requires
    showing that the nonlinear PSR formula with total SSV magnitude
    $|\Delta\SSV|^2 = |\Delta\SSV|_{\rm scalar}^2 +
    |(\SSV_{\rm net})_\phi|^2$ generates the exact functions
    $\Sigma = r^2 + a^2\cos^2\theta$ and $\Delta = r^2 - 2Mr + a^2$
    at all orders in $a/r$.  This requires the CPP field equation
    (companion~8, Proposition~4.2) evaluated in the rotating case.

  \item \textbf{Kerr--Newman extension.}
    Adding charge polarity SSV (companion~2~\cite{c2}) to the total
    SSV magnitude gives the Kerr-Newman metric for a charged, rotating
    black hole.  The mechanism is established; the explicit derivation
    is deferred.

  \item \textbf{Gravitational wave echoes from Kerr.}
    The GW echo calculation of companion~9~\cite{c9} used the
    Schwarzschild potential barrier.  For Kerr, the barrier depends
    on the spin parameter $a$ and the angular momentum $\ell$ of the
    perturbation mode.  The echo delay acquires a spin-dependent
    correction at order $(a/r_S)^2$.

  \item \textbf{Superradiance.}
    A rotating black hole can amplify waves that enter the ergosphere
    with frequencies below the superradiant threshold
    $\omega < m\Omega_+$ where $\Omega_+ = ac/(r_+^2 + a^2)$ is the
    horizon angular velocity and $m$ is the azimuthal mode number.
    The CPP mechanism for superradiance — azimuthal SSV coupling to
    the wave mode — is identified but not yet quantified.
\end{enumerate}

%======================================================================
\section{Conclusion}
\label{sec:conclusion}
%======================================================================

\begin{enumerate}
  \item \textbf{Azimuthal SSV source} (exact, §\ref{sec:source}):
    a rotating mass $J$ sources $k(\SSV_{\rm net})_\phi =
    GJ\sin^2\theta/(c^2 r^3)$ via the 12-edge selection rule.
    The $r^{-3}$ falloff distinguishes the vector source (angular
    momentum) from the scalar source (mass, $r^{-2}$ falloff).

  \item \textbf{Lense-Thirring frame-dragging} (exact, §\ref{sec:lense_thirring}):
    $g_{t\phi} = -2GJ\sin^2\theta/(c^2 r)$, derived with no free
    parameters and confirmed by Gravity Probe~B to $0.3\%$.

  \item \textbf{Kerr bound} (rigorous, Theorem~\ref{thm:kerr_bound}):
    $J \leq GM^2/c$ follows from SSV subluminality ($v_\phi \leq c$
    at the outer horizon), derived from the Absolute Moment postulate.
    Cosmic censorship is a theorem in CPP, not a conjecture.

  \item \textbf{Ergosphere} (rigorous, Theorem~\ref{thm:ergosphere}):
    the region where azimuthal SSV drag velocity equals $c$, bounded
    by $r_{\rm ergo} = M + \sqrt{M^2 - a^2\cos^2\theta}$.

  \item \textbf{Full Boyer-Lindquist Kerr metric}
    (Proposition~\ref{prop:kerr_full}): consistent with the CPP
    mechanism at all orders; the exact all-orders derivation from the
    CPP field equation is Open Problem~\ref{op:allorders}.

  \item \textbf{Planck-core ring} (from companion~8): CP Exclusion
    prevents the ring singularity; replaced by a Planck-density ring
    at the scale $\sim \lP$.

  \item \textbf{Schwarzschild limit exact} (Corollary~\ref{cor:schwarz_limit}):
    $J \to 0$ recovers companion~8 exactly.
\end{enumerate}

The CPP gravitational sector is now complete at the classical level:
Newtonian gravity (companion~5), weak-field GR and Eddington
factor-of-two (companion~7), exact Schwarzschild and Planck core
(companion~8), GW echoes (companion~9), Hawking radiation and Planck
remnant (companion~10), and Kerr frame-dragging and ergosphere
(companion~11).  All results follow from four postulates and the
600-cell geometry, with no free parameters beyond those fixed by the
Planck units.

%======================================================================
\section*{Acknowledgments}
%======================================================================

The azimuthal SSV source term and the Boyer-Lindquist consistency
argument were developed in collaboration with Grok (xAI).  The
rigorous derivation of the Kerr bound from SSV subluminality, the
ergosphere theorem, and the Schwarzschild limit corollary were
established on 18 March 2026.  The author thanks Claude (Anthropic)
for mathematical formulation, algebraic verification, and LaTeX
preparation.

%======================================================================
\begin{thebibliography}{99}

\bibitem{v16}
T.~L. Abshier,
``Mechanistic Derivation of Relativistic Effects via Space Stress
Vector (SSV) in the Dipole Sea,''
Version~16, 2026 (main paper).

\bibitem{c1}
T.~L. Abshier,
``The Absolute Moment Postulate,''
Version~3, 2026 (companion~1).

\bibitem{c2}
T.~L. Abshier,
``Microscopic Origin of the Dipole Sea Stiffness~$C$,''
Version~2, 2026 (companion~2).

\bibitem{c5}
T.~L. Abshier,
``Newtonian Gravity from SSV Shell Broadcast,''
Version~2, 2026 (companion~5).

\bibitem{c7}
T.~L. Abshier,
``Weak-Field General Relativity from the Lattice State Packet,''
Version~3, 2026 (companion~7).

\bibitem{c8}
T.~L. Abshier,
``Strong-Field General Relativity and the CPP Field Equation,''
Version~2, 2026 (companion~8).

\bibitem{c9}
T.~L. Abshier,
``Gravitational Wave Echoes from the Planck Core,''
Version~2, 2026 (companion~9).

\bibitem{c10}
T.~L. Abshier,
``Hawking Radiation and the Planck Remnant,''
Version~1, 2026 (companion~10).

\bibitem{kerr1963}
R.~P. Kerr,
``Gravitational Field of a Spinning Mass as an Example of
Algebraically Special Metrics,''
\textit{Phys.\ Rev.\ Lett.}, \textbf{11}, 237--238, 1963.

\bibitem{wald1984}
R.~M. Wald,
\textit{General Relativity},
University of Chicago Press, Chicago, 1984.

\bibitem{everitt2011}
C.~W.~F. Everitt \textit{et al.},
``Gravity Probe B: Final Results of a Space Experiment to Test
General Relativity,''
\textit{Phys.\ Rev.\ Lett.}, \textbf{106}, 221101, 2011.

\bibitem{ciufolini1995}
I.~Ciufolini and E.~C. Pavlis,
``A confirmation of the general relativistic prediction of the
Lense-Thirring precession,''
\textit{Nature}, \textbf{431}, 958--960, 2004.

\bibitem{lense1918}
J.~Lense and H.~Thirring,
``\"Uber den Einfluss der Eigenrotation der Zentralk\"orper auf die
Bewegung der Planeten und Monde nach der Einsteinschen
Gravitationstheorie,''
\textit{Phys.\ Z.}, \textbf{19}, 156--163, 1918.

\bibitem{penrose1969}
R.~Penrose,
``Gravitational Collapse: the Role of General Relativity,''
\textit{Riv.\ Nuovo Cimento}, \textbf{1}, 252--276, 1969.

\bibitem{codata}
NIST, CODATA 2018 Recommended Values of the Fundamental Physical
Constants, \url{https://physics.nist.gov/cuu/Constants/}, 2018.

\end{thebibliography}

\end{document}
